\section{Creation and Destruction of Blocks and Entities}\label{sec:creation-destruction}

As we have seen in \cref{sec:rendering}, we can use the topos of trees
$\toposoftrees$ to model streams, or more fittingly, finite approximations
thereof, which is done by modelling the passage of time, or of computation
steps, on the level of types.

We can in fact use this same construction to model \emph{ticks} (as they are
called in Minecraft), where a tick is a discrete step in time updating the
state. Consider again a Mathcraft
sheaf $W$ on $(V, \tau)$. We can lift it into the topos of trees by the
diagonal functor $\Delta$, yielding
\[
	\Delta W \in \functorcat{\omega^{\mathrm{op}}}{\functorcat{\tau^{\mathrm{op}}}{\Set}},
	\qquad
	\Delta W(k) = W \quad \text{for all } k \in \NN,
\]
where all restriction maps are identities.
\todo[inline]{Maybe instead this should be
	$\functorcat{\omega^{\mathrm{op}}}{\functorcat{\tau^{\mathrm{op}}}{\Set}}$?}

The interesting part, however, is the case in which we do not only have
identity maps. Let $W(0) = W$, and we will define the rest coinductively. We
wish to have a map $f \colon W(k+1) \to W(k)$ for each event $f$ that may
happen in $W$, like adding or removing blocks or entities, or mutating their
position.

% TODO
